\documentclass[10pt]{article}

\usepackage[utf8]{inputenc}

\usepackage[T1]{fontenc}

\usepackage{amsmath}

\usepackage{amsfonts}

\usepackage{amssymb}

\usepackage[version=4]{mhchem}

\usepackage{stmaryrd}

\usepackage{graphicx}

\usepackage[export]{adjustbox}

\graphicspath{ {./images/} }



\begin{document}

аннулятор порождается идемпотентом. Примерами бэровскіх колец служат кольцо эндоморфизмов векторного пространства над телом и кольцо ограниченных операторов гильбертова пространства. Бәровское кольцо наз. а б е л е в ы м, если все его идемпотенты центральны, и коне чны м (по Д е д е к и н д у), если $x y=1$ влечет за собой $y x=1$. Идемпотент $e$ бэровского кольца $R$ наз. а б е л е в м (к о н е ч ы м), если кольцо $e R e$ абелево (конечно по Дедекинду). Различаются следующие типы бэровских колец: $\mathrm{I}_{\text {fin }}$ - конечные кольца, содержащие абелев идемпотент, не принадлежащий никакому собственному прямому слагаемому; $\mathrm{I}_{\mathrm{inf}}-$ бесконечные по Дедекинду (т. е. не содержащие ненулевых конечных центральных идемпотентов) кольца, содержащие абелев идемпотент, не принадлежащий никакому собственному прямому слагаемому; $\mathrm{II}_{\mathrm{fin}}$ (или $I_{1}$ ) - конечные по Дедекинду кольца без ненулевых абелевых идемпотентов, содержащие конечный идемпотент, не принадлежащий никакому собственному прямому слагаемому; II inf - бесконечные по Дедекинду кольца с условием, указанным в II $\mathrm{I}_{\mathrm{fin}}$; III - кольца без ненулевых конечных идемпотентов. Каждое бәровское кольцо единственным способом разлагается в прямую сумму колец перечисленных типов (см. [9]) .



P. к. были введены для координатизации непрерывных геометрий, бирегулярные - в связи с исследованием функциональных представлений колец, бэровские (и риккатовы) - при исследовании колец операторов.



Рассматривались неассоциативные Р. к. См. также *-регуларное кальца, риккартово кольцо.



Jum.: [1] Б у р б а к и Н., Коммутативная алгебра, пер. с франц., М., 1971; [2] Д а у н с Д Ж., Г оф м а н К., «Математика", 1968 , т. 12, ㅅ 4, с. 3-24; [3] Л а м б е к И., Кольца и модули, пер. с англ., М., 1971; [4] С к о р н я к о в Ј. А., Дедекиндовы структуры с дополнениями и регулярные кольца, М., 1961; [5] Ф е й с К., Алгебра: кольца, модули и категории, пер. с



\begin{center}

\includegraphics[max width=\textwidth]{2023_07_08_59f942194af875d326a3g-1}

\end{center}



\includegraphics[max width=\textwidth, center]{2023_07_08_59f942194af875d326a3g-1(1)}

a r l K. R., Von Neumann regular rings, L. - [a. o.], 1979; [9] [10] N e u m a n n'J., Continuous geometry, Princeton,' 1960



*-РЕГУ ЛЯ РНОЕ Л. А. Скорнлков. допускающее инволюционный антиавтоморфизм $\alpha \mapsto \alpha^{*}$ такої, что $\alpha \alpha^{*}=0$ влечет за собой $\alpha=0$. Идемпотент $e$ из *-Р. к. наз. п р о е к ц и е й, если $e^{*}=e$. Каждый левый (правый) идеал *-Р.к. порождается однозначно определенной проекцией. Поәтому можно говорить о решетке проекций *-Р. к. Если әта решетка полна, то она является непрерывной геометрией. Дедекиндова решетка с дополнениями, обладающая однородным базисом $a_{1}, \ldots, a_{n}$, где $n \geqslant 4$ (см. Резулярное кольцо), является решеткой с ортодополнениями тогда и только тогда, когда она изоморфна решетке проекций нек-рого *-P. к.



Лum.: [1] С к о р н я к о в Л. А., Дедекиндовы структуры с дополнениями и регулярные кольца, М., 1961; [2] В е r b e r in a n S. K., Baer *-rings, B. - [a. o.], 1972; [3] K a p l a n sk y I., Rings of operators, N. Y.- Amst., 1968.



РЕГУ ЛЯ РНОЕ ПРЕДСТАВЛЕНИЕ - 1) Р. П. (Л ев о е) а л г е б р ы $A$ - линейное представление $L$ алгебры $A$ в векторном пространстве $E=A$, определяемое формулой $L(a) b=a b$ для всех $a, b \in A$; аналогично, формула $R(a) b=b a, a, b \in A$, определяет (анти)-представление алгебры $A$ в пространстве $E=A$, наз. (иравым) Р. п. $A$. Если $A$ - топологич. алгебра (с умножением, непрерывным по совокупности переменных), то $L$ и $R$ - непрерывные представлөния. Если $A$ алгебра с единицей или полупростая алгебра, то все ее Р. П.- точные.



\begin{enumerate}

  \setcounter{enumi}{1}

  \item Р. п. (п р а в о е) г р у п п ы $G$ - линейное представление $R$ группы $G$ в пространстве $E$ комплекснозначных функций на $G$, определенное формулой

\end{enumerate}



\[

(R(g) f)\left(g_{1}\right)=f\left(g_{1} g\right), g, g_{1} \in G, f \in E,

\]



шричем пространство $E$ разделяет точки групшы $G$ и облладает теџ свойствоџ, что функция $g_{1} \rightarrow f\left(g_{1} g\right), g_{1} \in G$, принадлежит пространству $E$ для всех $f \in E, g \in G$. Аналогично, формула



$(L(g) f)\left(g_{1}\right)=f\left(g^{-1} g_{1}\right), g, g_{1} \in G, f \in E$,



огределяет (левое) Р. п. группы $G$ в пространстве $E$, если функция $g_{1} \rightarrow f\left(g^{-1} g_{1}\right), \quad g_{1} \in G, \quad$ принадлежит $E$ для всех $g \in G, f \in E$. Если $G$ - топологич. группа, то в качестве пространства $E$ часто рассматриваются пространства непрерывных функций на $G$. Если $G$ - локально компактная группа, то (правым) Р. п. группы $G$ наз. (правое) Р. п. группы $G$ в пространстве $L^{2}(G)$, построенном по правоинвариантной мере Хаара на $G$; Р. п. локально компактной группы является ее непрерывным унитарным представлением, причем левое и правое Р. п. унитарно эквивалентны. А. И. Штерн



РЕГУ ЛЯ РНОЕ ПРОСТОЕ ЧИСЛО - простое нечетное число $p$, для к-рого число классов идеалов кругового поля $R\left(e^{2 \pi \imath / p}\right)$ не делится на $p$. Все остальные простые нечетные числа наз. и р р е г у л я р ны м и (см. Иррегулярное простое число). О. А. Иванова.



РЕГУЛЯРНОЕ ПРОСТРАНСТВО - топологИческое пространство, в к-ром для каждой точки $x$ и каждого не содержащего ее замкнутого множества $A$ найдутся непересекающиеся множества $U$ и $V$ такие, что $x \in U$ и $A \in V$. Регулярными являются все вполне регулярные пространства и, в частности, все метрические пространства.



Если в Р. п. все одноточечные подмножества замкнуты (а әто выполняется не всегда!), то оно наз. $T_{3}$-П р о с т р а в с т в о м. Не всякое Р. п. вполне регулярно: существует бесконечное $T_{3}$-пространство, на к-ром каждая непрерывная действительная функция постоянна Тем более, не каждое Р. п. является нормальным пространством. Однако если пространство регулярно и из каждого его открытого покрытия можно выделить счетное подпокрытие, то оно нормально. Пространство со счетной базой метризуемо в том и только в том случае, если оно является $T_{3}$-пространством. Регулярность наследуется любыми пространствами и мультипликативна.



Лит.: [1] $\mathrm{K}$ е л л и Д ж., Общая топология, пер. с англ., 2 изд., М., 1981; [2] А р х а н г е ль с к и й А В., П о н о м а p е в В. И., Основы общей топологии в задачах и упражнениях,



РЕГУ ЛЯ РНОЕ СОБЫТИЕ - множество слов конечного алфавита, к-рое на алгебраич. языке может быть задано с использованием выражений специального віда-рег у ля рных вы ра жени й. Пусть $A-$ конечный алфовит и $\mathrm{U}, \circ, *-$ символы операций, наз. объе динением, конкатен ацие й и и т е р а ц и е й соответственно. Р е г ул я р п ы в ы р а ж е и я в алфавите $A$ задаются индуктивно: 1) каждая буква из алфавита $A$ есть регулярное выражение, 2) если $R_{1}, R_{2}$ и $R$ - регулярные выражения, то $\left(R_{1} \cup R_{2}\right),\left(R_{1} \circ R_{2}\right)$ и $R^{*}$ суть также регулярные выражения. Язык регулярных выражений интерпретируется следующим образом.



Пусть $A *$ - множество всех слов в алфавите $A$, $\mathfrak{H}_{1} \subseteq A^{*}, \mathfrak{H}_{2} \subseteq A^{*}$. Символ любой буквы из $A$ понимается как множество, состоящее из одной буквы, а $\mathfrak{A}_{1} \cup \mathfrak{A}_{2}$ - как обычное теоретико-множественное объединение. Множество $\mathfrak{H}_{1} \circ \mathfrak{H}_{2}$ состоит из всех слов, к-рые представимы в виде $\alpha_{1} \alpha_{2}$, где $\alpha_{1} \in \mathfrak{A}_{1}, \alpha_{2} \in \mathfrak{A}_{2}$, причем если $\mathfrak{H}_{1}$ и $\mathfrak{H}_{2}$ содержат пустое слово $\Lambda$, то $\Lambda \alpha_{2}=\alpha_{2}, \alpha_{1} \Lambda=\alpha_{1}, \Lambda \Lambda=\Lambda$. Пусть $\mathfrak{A} \subseteq A^{*}$ и для любого $n \geqslant 2$ имеет место обозначение $\mathfrak{A} n=\mathfrak{P}$ и . . о $\mathfrak{A}$ ( $n$ раз). Тогда $\mathfrak{Y} *$ совпадает с $\mathfrak{A} \bigcup \mathfrak{A}^{2} \bigcup \mathfrak{P}^{3} \bigcup \ldots$, т. ө. если $\alpha \in \mathfrak{A} *$, то существует $m \geqslant 1$ такое, что $\alpha$ представимо в виде $\alpha_{1} \alpha_{2} \ldots \alpha_{m}$, где $\alpha_{i} \in \mathfrak{H}$ для любого $i=$ $=1, \ldots ., m$. Таким образом, множество слов конечного алфавита является ре гуля яны м с обы тием





\end{document}